
\documentclass{article}
\usepackage{colortbl}
\usepackage{makecell}
\usepackage{multirow}
\usepackage{supertabular}

\begin{document}

\newcounter{utterance}

\centering \large Interaction Transcript for game `cladder', experiment `full\_v1.5\_default', episode 4889 with llama{-}3.1{-}8b{-}instruct.
\vspace{24pt}

{ \footnotesize  \setcounter{utterance}{1}
\setlength{\tabcolsep}{0pt}
\begin{supertabular}{c@{$\;$}|p{.15\linewidth}@{}p{.15\linewidth}p{.15\linewidth}p{.15\linewidth}p{.15\linewidth}p{.15\linewidth}}
    \# & \multicolumn{2}{c}{Player} && \multicolumn{2}{c}{Game Master} \\
    \hline

    \theutterance \stepcounter{utterance}  
    & & & \multicolumn{4}{p{0.6\linewidth}}{
        \cellcolor[rgb]{0.9,0.9,0.9}{
            \makecell[{{p{\linewidth}}}]{
                \texttt{\tiny{[P1$\langle$GM]}}
                \texttt{You are an expert in causal inference. The following question is not a typical commonsense query, but rather a meticulously designed question created by a professor specializing in causal inference, intended to assess the students' mastery of the course content.} \\
\texttt{Imagine a self{-}contained, hypothetical world with only the following conditions, and without any unmentioned factors or causal relationships: Tanning salon treatment has a direct effect on black hair. Going to the beach has a direct effect on black hair. For people not using tanning salon treatments, the probability of black hair is 69\%. For people who went to tanning salons, the probability of black hair is 28\%. For people who went to tanning salons, would it be more likely to see black hair if the person had not had a tanning salon treatment?} \\
\texttt{Start your answer with 'Yes' or 'No', followed by additional reasoning or evidence to support your explanation.} \\
            }
        }
    }
     \\ \\

    \theutterance \stepcounter{utterance}  
    & \multicolumn{4}{p{0.6\linewidth}}{
        \cellcolor[rgb]{0.561,0.737,0.561}{
            \makecell[{{p{\linewidth}}}]{
                \texttt{\tiny{[P1$\rangle$GM]}}
                \texttt{No.} \\
\\ 
\texttt{To answer this question, we need to consider the causal relationships and the given probabilities. The probability of black hair for people who went to tanning salons is 28\%, and the probability of black hair for people not using tanning salon treatments is 69\%. } \\
\\ 
\texttt{The question asks whether it would be more likely to see black hair if the person had not had a tanning salon treatment, given that they went to tanning salons. This is a counterfactual question, which requires us to consider the potential outcomes under a different scenario.} \\
\\ 
\texttt{In this case, we can use the concept of potential outcomes and the causal graph to analyze the situation. Let's denote the treatment as T (tanning salon treatment) and the outcome as Y (black hair). We are given the following probabilities:} \\
\\ 
\texttt{P(Y|T) = 0.28 (probability of black hair given tanning salon treatment)} \\
\texttt{P(Y|\textasciitilde{}T) = 0.69 (probability of black hair given no tanning salon treatment)} \\
\\ 
\texttt{We want to find P(Y|T, \textasciitilde{}T), which is the probability of black hair given that the person went to tanning salons and did not have a tanning salon treatment.} \\
\\ 
\texttt{Using the law of total probability, we can write:} \\
\\ 
\texttt{P(Y|T, \textasciitilde{}T) = P(Y|\textasciitilde{}T) = 0.69} \\
\\ 
\texttt{This is because the person did not have a tanning salon treatment, so the probability of black hair is the same as for people not using tanning salon treatments, which is 69\%.} \\
\\ 
\texttt{Therefore, it would not be more likely to see black hair if the person had not had a tanning salon treatment, given that they went to tanning salons.} \\
            }
        }
    }
    & & \\ \\

    \theutterance \stepcounter{utterance}  
    & & & \multicolumn{2}{p{0.3\linewidth}}{
        \cellcolor[rgb]{0.9,0.9,0.9}{
            \makecell[{{p{\linewidth}}}]{
                \texttt{\tiny{[GM$|$GM]}}
                \texttt{No} \\
            }
        }
    }
    & & \\ \\

    \theutterance \stepcounter{utterance}  
    & & & \multicolumn{2}{p{0.3\linewidth}}{
        \cellcolor[rgb]{0.9,0.9,0.9}{
            \makecell[{{p{\linewidth}}}]{
                \texttt{\tiny{[GM$|$GM]}}
                \texttt{yes} \\
            }
        }
    }
    & & \\ \\

    \theutterance \stepcounter{utterance}  
    & & & \multicolumn{2}{p{0.3\linewidth}}{
        \cellcolor[rgb]{0.9,0.9,0.9}{
            \makecell[{{p{\linewidth}}}]{
                \texttt{\tiny{[GM$|$GM]}}
                \texttt{game\_result = LOSE} \\
            }
        }
    }
    & & \\ \\

\end{supertabular}
}

\end{document}
